\chapter{Tokens, Temperatura e Controlo de Loucura}
\label{ch:tokens}

Uma peça de roupa pode chegar com a etiqueta «lã», um código e um desenho minúsculo de uma bacia. Para nós, cada um tem um significado. Para contar o que cabe na entrada de um modelo de texto, precisamos de conhecer a forma como o tokenizador representa esses sinais.

\section{Como o texto se torna uma sequência de IDs}
O tokenizador converte o texto numa sequência de identificadores do seu vocabulário. Um token pode representar uma palavra inteira, parte dela, pontuação ou uma sequência que inclui espaços. Usa o tokenizador que corresponde à versão e configuração do modelo, pois a segmentação pode variar.

\subsection*{Espaço e custos}
Quando o serviço dá um limite em tokens, conta com o texto e com a estrutura que a aplicação acrescenta, conforme a interface. Cinco páginas de código e cinco páginas de prosa podem ter contagens diferentes. A contagem serve para decidir o que cabe e, quando a cobrança usa tokens, para estimar parte do custo.

\subsection*{Dois textos que parecem iguais}
Em Python, podemos representar «é» por um ponto de código ou por um \texttt{e} seguido de um acento combinante. O seguinte cálculo mostra o número de pontos de código e de bytes em UTF-8 de cada representação.

\begin{lstlisting}
>>> [ (len(s), len(s.encode("utf-8")))
...   for s in ("\u00e9", "e\u0301") ]
[(1, 2), (2, 3)]
\end{lstlisting}

A contagem de tokens depende ainda do tokenizador e da normalização que este aplicar. Para contar palavras, letras ou sílabas numa saída, define a regra da tarefa e verifica segundo essa regra. Um limite de tokens controla outro tamanho.

\begin{promptlabbox}{Ver a tokenização com os teus próprios olhos}
\textbf{Objetivo e contexto:} Medir segmentação com um tokenizador identificado. Guarda os três textos completos e a versão usada.

Usa um tokenizador fornecido ou documentado para o modelo escolhido. Se a interface só mostrar a contagem de utilização, regista essa contagem e a diferença entre o pedido escrito e a chamada completa.

\medskip

Tarefa:

\begin{enumerate}[leftmargin=*]
\item Escolhe três frases: uma normal, uma com muitos acentos/abreviaturas, e uma com emojis.
\item Usa uma ferramenta oficial de tokenização para ver em quantos tokens cada frase é partida.
\item Repara como pequenas mudanças (retirar um emoji, mudar uma palavra) alteram o número de tokens.
\end{enumerate}

Depois disso, nunca mais vais olhar para “contar palavras” da mesma forma.

\textbf{Resultado esperado e validação:} Regista IDs ou segmentos e contagem de tokens. Compara tamanho, sem presumir que um acento ou emoji aumenta sempre a contagem.

\textbf{Falhas possíveis:} Usar tokenizador de outro modelo, confundir caracteres com tokens ou introduzir normalização sem a registar.
\end{promptlabbox}

\section{Parâmetros de \emph{sampling}: temperatura, top-k e top-p}
O modelo produz \emph{logits}, scores reais $z_i$ para tokens candidatos. Para $T>0$, uma forma comum de obter probabilidades é
\[
p_i(T)=\frac{\exp(z_i/T)}{\sum_j\exp(z_j/T)}.
\]
Reduzir $T$ concentra a distribuição nos scores maiores. Num exemplo com dois tokens e logits $\ln 2$ e $0$, o primeiro tem probabilidade $2/3$ a $T=1$. A $T=0{,}5$, tem $4/5$. A $T=2$, desce para cerca de $0{,}586$. Pensa na temperatura como controlo da distribuição de escolha. Se o token mais favorecido fizer parte de uma resposta errada, baixar $T$ dá-lhe mais peso.

\emph{Top-k} restringe a escolha aos $k$ candidatos de maior probabilidade. \emph{Top-p}, ou nucleus sampling, retém o menor prefixo da lista ordenada cuja probabilidade acumulada atinge pelo menos $p$; depois renormaliza~\citep{holtzman2020}. Empates e combinação de filtros dependem da implementação.

Com probabilidades ilustrativas $[0.60,0.25,0.10,0.05]$, top-p de $0.80$ retém os dois primeiros candidatos, cuja massa é $0.85$. Depois de renormalizar, as probabilidades são $12/17$ e $5/17$. Top-k igual a 1 corresponde a escolha gulosa no passo considerado. O limite $T\to0^+$ também concentra a escolha nos máximos; muitos serviços tratam o valor zero como opção especial. Não substituas essa convenção pela divisão por zero na fórmula.

\begin{warningboxenv}{Repetível não significa garantidamente idêntico}
Mesmo sem amostragem, uma aplicação pode mudar de versão, contexto, ferramenta ou implementação numérica. Temperatura zero não é contrato de reprodução de um serviço inteiro. Guarda a versão do modelo, configuração, entradas e resultados das ferramentas.
\end{warningboxenv}

Nem todos os modelos expõem estes controlos, e alguns modelos de raciocínio impõem restrições específicas. Consulta a documentação da versão que vais usar. Em experiências, muda um parâmetro de cada vez e mede correção, diversidade, custo e latência; não adotes «0.7» como constante da natureza.


\section{Orientar o estilo quando não controlas a API}

Nem sempre tens acesso direto aos parâmetros de \emph{sampling}. Se o chat não expõe temperatura ou top-p, podes descrever o tom e o formato que queres. Ainda assim, consegues influenciar o comportamento com instruções.

Alguns padrões úteis:

\begin{itemize}[leftmargin=*]
\item \textbf{Para pedir redação contida:}
\begin{itemize}[leftmargin=*]
\item “Responde de forma factual e direta.”
\item “Evita floreados estilísticos.”
\item “Se não tiveres a certeza, diz explicitamente que não sabes.”
\end{itemize}
\item \textbf{Para pedir exploração de ideias:}
\begin{itemize}[leftmargin=*]
\item “Sê criativo e explora possibilidades improváveis.”
\item “Dá-me várias ideias muito diferentes umas das outras.”
\end{itemize}
\item \textbf{Para controlar variação:}
\begin{itemize}[leftmargin=*]
\item “Mantém o estilo consistente com o exemplo seguinte.”
\item “Usa sempre o mesmo formato que já definimos acima.”
\end{itemize}
\end{itemize}

Estas instruções mudam a distribuição condicionada ao texto. Não alteram o parâmetro numérico de temperatura, não substituem top-p e não oferecem garantias equivalentes. «Responde de forma factual» é um objetivo; verificar os factos continua a ser outra operação.

\begin{promptlabbox}{Criar um seletor de “modo” só com texto}
\textbf{Objetivo e contexto:} Testar três instruções de estilo na mesma tarefa e com os mesmos dados.

Constrói um prompt de system com instruções deste género:

\medskip

\begin{lstlisting}
És um assistente de escrita com três modos:

[modo=PRECISO]

Estilo: direto, conciso, foco em factos.

Evita criatividade desnecessária.

Admite quando não souberes algo.

[modo=EXPLICATIVO]

Estilo: pedagógico, com exemplos.

Mantém rigor técnico, mas admite analogias.

[modo=CRIATIVO]

Estilo: exploratório, ideias fora da caixa.

Permite sugestões improváveis, desde que coerentes.

Sempre que o utilizador escrever algo como:
[MODO: PRECISO] ou [MODO: EXPLICATIVO] ou [MODO: CRIATIVO],
deves adaptar o teu comportamento a esse modo.
\end{lstlisting}

Depois, experimenta:

\begin{itemize}[leftmargin=*]
\item Fazer a mesma pergunta nos três modos.
\item Ver quão diferentes são as respostas.
\end{itemize}

\textbf{Resultado esperado e validação:} Compara concisão, exemplos e diversidade com critérios textuais. Regista também erros factuais; uma mudança de estilo não demonstra alteração da temperatura da API.

\textbf{Falhas possíveis:} Confundir criatividade verbal com qualidade ou deixar modos diferentes acumularem contexto.
\end{promptlabbox}

\section{Comprimento, penalizações e repetição}

Outro conjunto de parâmetros importante controla o quão longa e repetitiva é a resposta do modelo.

Nas APIs, encontras coisas como:

\begin{itemize}[leftmargin=*]
\item \texttt{max\_tokens} ou equivalente: limite de geração; o nome e a inclusão de tokens de raciocínio variam por API.
\item Penalização por frequência, quando disponível: pode depender de quantas vezes um token foi visto.
\item Penalização por presença, quando disponível: pode depender de o token já ter aparecido.
\item Confirma na documentação a fórmula e os tokens a que cada parâmetro se aplica.
\end{itemize}

Se não tens acesso direto a estes controlos, podes aproximar-te com linguagem natural:

\begin{itemize}[leftmargin=*]
\item Para limitar comprimento:
\begin{itemize}[leftmargin=*]
\item “Responde em no máximo 5 frases.”
\item “Dá-me apenas uma lista com 3 pontos.”
\end{itemize}
\item Para evitar repetição:
\begin{itemize}[leftmargin=*]
\item “Não repitas o enunciado.”
\item “Evita repetir as mesmas palavras; usa sinónimos.”
\end{itemize}
\end{itemize}

Claro que isto não é perfeito. O modelo pode ignorar-te ocasionalmente. Mas é muito melhor do que mandar perguntas vagas e esperar milagres.

\begin{warningboxenv}{O mito do “responde em exatamente N palavras”}
Pedir “exatamente 137 palavras” é pedir ao modelo que alinhe um pedido em “palavras” com uma realidade interna em “tokens”. Mesmo que acerte, não contes com isso como contrato. Em contextos de produção, usa limites de tokens na API e valida o tamanho do output com código.
\end{warningboxenv}

\section{Quando não usar LLMs}

Antes de avançarmos, convém dizer o óbvio que muitas equipas ignoram: há situações em que usar um LLM é simplesmente má ideia.

Casos típicos onde \emph{não} é a melhor ferramenta:

\begin{itemize}[leftmargin=*]
\item \textbf{Cálculos exatos e de alto risco}: impostos, dosagem de medicamentos, controlo de infraestruturas críticas.
\item \textbf{Verificações formais}: provas matemáticas, validação de protocolos criptográficos.
\item \textbf{Regras fixas e simples}: transformações deterministas (por exemplo, converter \texttt{YYYY-MM-DD} em \texttt{DD/MM/YYYY}).
\item \textbf{Dados ultra-sensíveis} quando não tens garantias fortes de privacidade.
\end{itemize}

Nestes casos, usa código normal, algoritmos tradicionais, bases de dados. O LLM pode ajudar a \emph{escrever} esse código, a explicar o problema ou a gerar testes — mas não deve ser o motor principal de decisão.

\begin{keyideabox}{LLM é martelo, mas nem tudo é prego}
Sempre que pensares “podíamos pôr um LLM aqui”, pergunta primeiro se não há uma função determinista simples que faça melhor, mais rápido e mais barato. Só depois envolve tokens.
\end{keyideabox}
